\section{Discussion}
\label{sec:discussion}

\subsection{Where the canopy drag is applied}

This study asked which modelled processes control transport in a shallow
vegetated lagoon, and the answer the ensemble gives is the seagrass canopy, by a
margin that leaves the other two treatments second order. The mechanism the
model offers for that is vertical rather than frictional, since introducing the
canopy removes momentum from the lower part of the water column without removing
a proportionate amount from the upper part, roughly doubling the ratio of
surface to bed speed. Such a profile is the signature of a canopy shear layer
rather than of a rougher bed, which acts through a logarithmic profile whose
surface-to-bed ratio is bounded well below the value observed here
\citep{Nepf2012}, and the bed shear stress confirms the distinction by its sign,
falling under the canopy where an equivalent bed roughness raises it over the
same cells. Since a surface drifter samples the top of that profile, a treatment
that reshapes the profile relocates the drifter even where it changes the
depth-averaged transport little.

That places the present result in a particular relation to the two earlier
models of this basin rather than in opposition to them.
\citet{Ingrassia2024} obtained a best-fit uniform Gauckler--Strickler
coefficient for the vegetated lagoon within a depth-averaged framework, where an
equivalent bed roughness is the appropriate representation of a canopy because
there is no vertical profile for the canopy to shape. \citet{DeMarchis2012}
worked in three dimensions and assigned Nikuradse roughness heights by seagrass
class, establishing the vertical structure of the wind- and tide-driven
circulation that motivates using a three-dimensional solver here at all. Both
answered the questions they set, and a roughness representation reproduces the
drag the meadow exerts. What has changed since is the availability of a
canopy formulation that resolves that drag over the vertical extent the plants
occupy, and the contribution of the present work is to quantify what that
capability buys in a basin of this depth. The canopy parameters used here were
taken from published measurement of the two species and were not adjusted
against any observation, so the improvement follows from where the drag is
applied rather than from how it was tuned.

\subsection{What each gauge position constrains}

Water level separated the configurations at AltaVilaEst and at neither of the
gauges near the inlets, and the three positions constrain different parts of the
problem. Agreement at the inlet gauges establishes that the offshore forcing,
the datum and the exchange through the two mouths are right, and in a micro-tidal
basin with two open connections those are what set the level there, so every
configuration reproduces it equally well. Agreement in the interior tests
something the outer pair cannot, because the level there is set by what the
water does on its way in, which is where the configurations differ. The two
roles are complementary, and a match in the interior would be difficult to
interpret without the exchange having been established first.

Together with the drifter result, the gauge result refines the familiar argument that water level
constrains transport weakly and that Lagrangian observations form a distinct
constraint rather than a confirmation \citep{Liu2011, Revelard2021}. Water level
is not inherently insensitive to the canopy, since the one gauge sited inside
the vegetated basin identified the same factor as the drifters and ranked the
configurations the same way. For a basin of this type the useful question is
therefore less whether to validate on water level than where to place the
instrument intended to discriminate, and the interior of the vegetated
compartment is a considerably more informative position for that purpose than
the inlet.

\subsection{What an ensemble without a canopy reports}
\label{sec:disc_synergy}

Taken on its own, the bare half of this design supports a strong statement about
process coupling. Wave coupling and bed mobility each reduce Lagrangian skill
when applied alone and raise it when applied together, with a large interaction
term, tight bootstrap intervals and a mechanism available in the literature,
since wave orbital motion sets the bed stress that mobilises sediment and
wave--tide coupling is known to govern the sediment dynamics of shallow
micro-tidal basins \citep{Carniello2011}.

That statement weakens considerably once the canopy is present, since in the
vegetated half the same interaction term falls by most of its value, both
contrasts lose their significance, and the spread of skill across the four
configurations contracts by a factor of three.

Compensation is the reading the rest of the evidence supports, in that without a
canopy the modelled profile is too uniform and a wave field acting on a bed that
can respond reproduces by another route part of the structure the meadow would
have produced, leaving nothing for that combination to reproduce once the meadow
is in place. This remains an inference from the design rather than a
demonstration of the mechanism, though it is consistent with the vertical
evidence, where the canopy moves the surface-to-bed ratio by an order of
magnitude more than wave coupling does.

What follows from this bears on how such an ensemble is read rather than on this
basin alone, since canopy drag is commonly left out of lagoon and estuary models
or folded into a bed roughness coefficient, and an ensemble built that way can
still return a well-defined and statistically clean interaction between the
processes it does carry. Here that interaction is confined to the configurations
without a canopy, so it describes those configurations rather than the system,
and nothing internal to the bare half of the design distinguishes the two
readings. What distinguished them was varying the canopy as a factor in its own
right and scoring the result against an observation type that no model decision
had used.

\subsection{Limitations}
\label{sec:disc_limits}

\textbf{Wave dissipation by the canopy.} The canopy does not reach the wave
model, because the coupling passes a bed roughness field to the spectral model
while the canopy enters the momentum equation of the flow, so mean orbital
velocity at the bed
changes by under 2.5\% when the canopy is introduced and not in an order that
follows canopy density. Seagrass meadows do attenuate waves, in this species and
at this scale \citep{Infantes2012, Prinos2010}, and shallow restricted
Mediterranean lagoons carry wave fields that respond to bed conditions
\citep{Morote2025}. The wave field in these simulations is therefore that of a
bare bed, and the wave contrasts reported here are conditional on that.
Representing canopy dissipation inside the spectral model is the corresponding
extension. It would be expected to change the wave contrasts rather than the canopy
ones, which are measured on the flow.

\textbf{Canopy reconfiguration.} The canopy bends under load rather than holding
a fixed height, and the response is dynamic, the velocity ratio between two runs
differing only in leaf stiffness moving from 1.005 in the slowest quintile of
the record to 0.807 in the fastest 5\%, where a fixed shortening would give a
constant ratio. A 13\% change in leaf density moves the velocity field by 20\%
at the upper end of the range, which is a strong sensitivity to a value taken
from a published table rather than measured in this meadow. The bending is
moreover mechanical, whereas the flume measurements of \citet{CiraoloFerreri2006} describe
reconfiguration empirically through a friction factor that varies by a factor of
twelve over the velocity range they tested. That empirical formulation is
available in the model only in its depth-averaged solver, so the two could not
be compared directly and the agreement between them was not tested.

\textbf{Bed mobility.} The morphodynamic factor is not calibrated, since the
critical shear stress for erosion has not been constrained against a repeated
bathymetric survey, and over nine days the modelled bed changes by up to 1.4~m
in isolated cells against a median of 0.05~m in a basin of 0.95~m mean depth,
which is not a plausible rate for this system. The factor is therefore a
sensitivity test rather than a prediction, and its bearing on the conclusions is
limited by the best-performing configuration of the eight carrying a fixed
bed.

\textbf{Window and restart segments.} The record covers two days of water level
and thirty-four hours of drifter observations within one summer season, and the
vegetated configurations are three-day restart segments. A control run
establishes that the restart accounts for none of the water level result or of
the drifter skill (Section~\ref{sec:res_wl}), but seasonal
behaviour, storm conditions and the interaction of a canopy with a mobile bed
over longer periods are outside what this design addresses.

\textbf{Initial condition.} Started from the stratified condition instead of the
uniform offshore column of the ensemble, the canopy contrast on the uncoupled
fixed bed is 0.120 rather than 0.195, since the uniform start lowers the skill of the bare member
by 0.080 and leaves the vegetated member within 0.005
(Section~\ref{sec:res_numerics}). The contrast remains clear of zero under
either start, and water level responds only by a shift of its mean, but the
magnitudes reported here, those of the bare arm in particular, are conditional
on the uniform start. The solution also depends on the host and the
decomposition at the Courant limit of the ensemble, by 0.034 in skill, which all
members share, while the horizontal and vertical resolution do not limit the
reported quantities.

\textbf{\emph{Cymodocea nodosa} parameters.} The \emph{C.\ nodosa} parameters
are loosely constrained, coming from a flume study of plants collected
elsewhere, since no measurement of that species in this lagoon was available.
The species carries about 1.5\% of the canopy frontal area in the model, which
bounds the consequence, and the drag is carried almost entirely by the
\emph{P.\ oceanica} column.
